\documentclass{article}
\usepackage{amsmath, amssymb, amsthm}

% Define standard theorem environments
\newtheorem{theorem}{Theorem}
\newtheorem{lemma}{Lemma}
\newtheorem{definition}{Definition}
\newtheorem{axiom}{Axiom}

\begin{document}
\section*{Theoretical Framework: From Local Trajectories to Global Latent Space}

Let $\mathcal{X}$ denote the universal space of household observations. We partition
this space into a finite set of discrete, non-overlapping geographies $g \in G$. To
accommodate continuous dynamics, let time $\mathcal{T} \subset \mathbb{R}$ be a
continuous interval, and $\mathcal{S} \subset \mathbb{R}^n$ denote the cross-sectional
state space of household attributes.

The full set of observations for a specific geography is defined as the product space
$X_g = \mathcal{T} \times \mathcal{S}$. We equip $X_g$ with the standard Euclidean
subspace topology. The universal space is the disjoint union
$\mathcal{X} = \bigsqcup_{g \in G} X_g$, equipped with the disjoint union topology
(meaning a subset is open in $\mathcal{X}$ if and only if its intersection with every
geography $X_g$ is open).

Let $I$ be a set of indexed households. A specific household $i \in I$'s trajectory
through time is given by a transport function $f^{g,i}_{t,k} : \mathcal{S} \to \mathcal{S}$,
which maps household $i$'s realized state at time $t$ to its own realized state at time
$k$ within the same geography.

We introduce an economic ordering $\succsim$ representing real wealth. Rather than
assuming universal global comparability outright---which is fraught with spatial and
intertemporal measurement issues---we impose comparability only along strictly defined,
localized empirical dimensions. By explicitly leaving certain pairs incomparable,
$\succsim$ constitutes an \textit{incomplete} empirical preference relation.

\textbf{Axiom 1} (Restricted Domain of Comparability). \textit{Let $\succsim$ be an
empirical binary relation representing real wealth. $\succsim$ is strictly defined only
under the following localized conditions:}
\begin{itemize}
    \item \textbf{Cross-Sectional Comparability:} Within any geography $g$ and a fixed
    time $t \in \mathcal{T}$, any two household states $x = (t, s_x)$ and
    $y = (t, s_y)$ in $X_g$ are directly comparable.
    \item \textbf{Trajectory Comparability:} For a specific household $i$'s trajectory,
    its state at time $t$, $x = (t, s) \in X_g$, and its transported state at time $k$,
    $x' = (k, f^{g,i}_{t,k}(s)) \in X_g$, are directly comparable.
    \item \textbf{Orthogonal Incomparability (Intertemporal Cross-Household):} The
    relation cannot be evaluated across both dimensions simultaneously in the empirical
    data. For distinct households $i \neq j$ observed at distinct times $t \neq k$,
    direct comparison is strictly undefined ($x_{i,t} \parallel x_{j,k}$).
    \item \textbf{Spatial Incomparability:} For any two observations in different
    geographies, $x \in X_g$ and $y \in X_h$ (where $g \neq h$), the relation is
    strictly undefined ($x \parallel y$).
\end{itemize}

\textbf{Axiom 2} (Acyclicity and Continuous Transitive Closure). \textit{Transitivity
of the empirical relation $\succsim$ holds strictly within the defined domains in
Axiom~1, but does not explicitly chain across orthogonal domains. However, the global
relation $\succsim$ on $X_g$ is strictly acyclic. Furthermore, we assume the transitive
closure of this relation, denoted $\succsim^*$, forms a continuous partial preorder on
$X_g$, meaning its upper and lower contour sets are closed under the standard product
topology of $X_g$.}

Because we embrace the natural incompleteness of the empirical data, we do not require
the topological space to be perfectly ordered. Instead, we demonstrate that a global
machine learning architecture can validly map this disjoint and partially ordered space
into a single continuous scalar dimension, preserving the known data while bridging the
gaps via continuous spatial projection.

\begin{lemma}[Local Continuous Representation]
For each geography $g \in G$, there exists a continuous function
$u_g : X_g \to \mathbb{R}$ such that for all $x, y \in X_g$:
$$x \succ^* y \implies u_g(x) > u_g(y)
  \quad\text{and}\quad
  x \sim^* y \implies u_g(x) = u_g(y),$$
where $\succ^*$ and $\sim^*$ denote the strict and symmetric parts of the transitive
closure $\succsim^*$, respectively.
\end{lemma}

\begin{proof}
Fix an arbitrary geography $g \in G$. We verify the two hypotheses required to apply
the continuous representation theorem for partial preorders (Herden, 1989; Ok, 2002).

\medskip
\noindent\textbf{Step 1: Second-countability of $X_g$.}
By definition, $X_g = \mathcal{T} \times \mathcal{S}$ where $\mathcal{T} \subset \mathbb{R}$
and $\mathcal{S} \subset \mathbb{R}^n$, so $X_g$ embeds as a subspace of $\mathbb{R}^{1+n}$.
The Euclidean space $\mathbb{R}^{1+n}$ is second-countable, admitting a countable base of
open balls with rational centers and radii. Since second-countability is hereditary under
subspaces, $X_g$ is second-countable under the Euclidean subspace topology.

\medskip
\noindent\textbf{Step 2: $\succsim^*$ is a continuous partial preorder on $X_g$.}
We verify the three required properties.

\textit{Reflexivity.} The base relation $\succsim$ is an economic ordering representing
real wealth; any household state is at least as wealthy as itself, so $x \succsim x$ for
all $x \in X_g$. This base reflexivity is inherited by $\succsim^*$.

\textit{Transitivity.} Transitivity holds by the construction of the transitive closure.

\textit{Continuity.} Axiom~2 stipulates that all upper contour sets
$U(x) = \{y \in X_g : y \succsim^* x\}$ and lower contour sets
$L(x) = \{y \in X_g : x \succsim^* y\}$ are closed in the product topology of $X_g$.
This is precisely the definition of a \textit{continuous} partial preorder.

Together, these three properties confirm that $(X_g, \succsim^*)$ is a continuous
partial preorder. Note that completeness is not required.

\medskip
\noindent\textbf{Step 3: Existence of $u_g$.}
A continuous partial preorder with closed contour sets on a second-countable topological
space admits a continuous, order-preserving scalar representation (Herden, 1989,
Theorem~1; Ok, 2002, Theorem~1). Applying this result to $(X_g, \succsim^*)$ yields the
existence of a continuous function $u_g : X_g \to \mathbb{R}$ satisfying the stated
order-preservation conditions. Since $\succsim^*$ is the transitive closure of the base
empirical relation $\succsim$, every comparison defined by $\succsim$ is preserved by
$\succsim^*$, and hence by $u_g$.
\end{proof}

\begin{theorem}[Global Latent Representation]
If the incomplete relation $\succsim$ on the space $\mathcal{X}$ satisfies Axioms~1
and~2, there exists a global continuous scalar function
$\hat{U} : \mathcal{X} \to \mathbb{R}$ providing a forward representation, such that
for any empirically defined comparable pair $x, y \in X_g$:
$$ x \succ y \implies \hat{U}(x) > \hat{U}(y)
   \quad\text{and}\quad
   x \sim y \implies \hat{U}(x) = \hat{U}(y). $$
\end{theorem}

\begin{proof}
By Lemma~1, for each geography $g \in G$ there exists a continuous, order-preserving
local function $u_g : X_g \to \mathbb{R}$ representing $\succsim^*$ on $X_g$. Define
$\hat{U} : \mathcal{X} \to \mathbb{R}$ piecewise by
$$ \hat{U}(x) = u_g(x) \quad \text{where } x \in X_g. $$

\noindent\textbf{Global continuity.}
The universal space $\mathcal{X} = \bigsqcup_{g \in G} X_g$ carries the disjoint union
topology, under which a function is continuous if and only if its restriction to each
component $X_g$ is continuous. Since every $u_g$ is continuous by Lemma~1, $\hat{U}$
is globally continuous on $\mathcal{X}$.

\noindent\textbf{Order-preservation.}
Let $x, y \in X_g$ be any empirically comparable pair. Since $\succsim$ is defined
within a single geography (Axiom~1) and $u_g$ represents $\succsim^*$ on $X_g$, the
stated forward implications hold immediately.

\noindent\textbf{No violation of incomparability.}
By Axiom~1 (Spatial Incomparability), $\succsim$ defines no comparisons between $X_g$
and $X_h$ for $g \neq h$. Because $\hat{U}$ is only claimed to preserve comparisons
defined within a single $X_g$, projecting disjoint geographies into the shared codomain
$\mathbb{R}$ introduces no contradiction with any established economic relation.
\hfill$\square$
\end{proof}

\textbf{Remark on AI Latent Space and Residual Projections:}
Axiom~1 explicitly leaves the intertemporal relationship between distinct households,
as well as cross-geography comparisons, formally undefined in the empirical data.
However, a neural network maps these disjoint and partially ordered topologies into a
shared, totally ordered latent space ($\mathbb{R}$).

Because $\mathbb{R}$ is totally ordered, for any two mathematically incomparable
observations $x, y \in \mathcal{X}$, the model will inherently output scalar values such
that either $\hat{U}(x) \ge \hat{U}(y)$ or $\hat{U}(y) \ge \hat{U}(x)$. These forced
comparisons are \textbf{residual mathematical artifacts} of the latent projection (which
is why Theorem~1 only claims a one-way, forward implication). They do not come from
unfounded empirical assumptions forced onto the raw data, but rather emerge naturally
and interpretably as the machine learning architecture interpolates across the shared
continuous spatial and attribute features.

Furthermore, this topological formulation maintains a critical advantage for empirical
measurement: it does not require the household physical transport function ($f$) to be
continuous. Empirical data can be noisy and households can experience discontinuous
shocks. As long as the underlying economic valuation of those states (Axiom~2) is
continuous, a smooth latent space mapping remains mathematically valid and strictly
preserves the true local intertemporal economic dynamics established in the data.

\end{document}